% bdlt-units: cm
% deck.tex -- the guide deck of BDLT PDF Presenter: the PDF that opens first (public/sample.pdf).
%
%   python3 <bdlt-beamer-media>/tools/bdlt-deck.py build deck.tex          the presenter's guide
%   python3 <bdlt-beamer-media>/tools/bdlt-deck.py build deck-viewer.tex   the viewer's guide (students)
%   ./build.sh                                                             both, and copies them to public/
%
% One source for both editions: \ifguideviewer is true in deck-viewer.tex, which leaves out what
% only a lecturer needs. Every picture is a vector drawing (guide-icons.tex, made by
% make-icons.py; cover.pdf; the posters of the two live figures), so the PDF stays small:
% every visitor of the presenter downloads it. For the same reason both live figures come from
% one embedded file, media/figures/guide-figures.html (the key x-figure chooses), so that the
% PDF carries Source Sans and D3 once; each figure has its own source in media/<name>/.
\documentclass{beamer}
\usetheme[media]{BDLT26}
\bdltmediasetup{mediapath={media/}}
\graphicspath{{media/}}
\usetikzlibrary{svg.path}
% poster.tex -- the poster of a live figure as a TikZ drawing, laid under the figure's place.
%
%   \guideposter{<name>-poster.tex}{W}{H}\bdltwidget[..., poster=blank-poster.pdf]{<figure>.html}
%
% <name>-poster.tex is written by make-poster-tikz.mjs from the figure itself (W x H CSS px, the
% AREA of the figure's build.sh). It is drawn where the theme would put a poster image: fitted into
% the picture area of the layout and centred. The theme still wants a poster file, so the figure
% names blank-poster.pdf, an empty page. Why not a poster image: text in an image brings its own
% letters with it (45 kB for the two posters of this deck); text set by TeX uses the deck's.
% Use it inside a picture area of the theme (\bdltfullimage, \bdltarea, ...). Needs TikZ with svg.path.
\newlength\posterlen
\newcommand\guideposter[3]{%
  \posterlen=\dimexpr\bdltareawidth/#2\relax
  \ifdim #3\posterlen>\bdltareaheight \posterlen=\dimexpr\bdltareaheight/#3\relax\fi
  \makebox[0pt][l]{\hskip\dimexpr(\bdltareawidth-#2\posterlen)/2\relax
    \raisebox{\dimexpr(\bdltareaheight-#3\posterlen)/2\relax}{\input{#1}}}}
   % \guideposter: the posters of the live figures as TikZ drawings

\newif\ifguideviewer
\ifdefined\guideviewer \guideviewertrue \fi

\title{\ifguideviewer BDLT PDF Viewer\else BDLT PDF Presenter\fi}
\subtitle{\ifguideviewer How to read, annotate and save\else How to present, annotate and save\fi}
\author{Blockchain \& DLT Research Group, UZH}
% the release of the presenter this guide describes: build.sh reads it from public/index.html
\IfFileExists{guide-version.tex}{%% guide-version.tex -- written by build.sh from public/index.html; do not edit.
\def\guideversion{1.6.2}
}{\def\guideversion{}}
\date[Guide]{\ifx\guideversion\empty Guide\else Guide to version \guideversion\fi}
\titlegraphic{cover}

% ---------------------------------------------------------------- drawing helpers (lengths in cm)
\definecolor{guidetile}{HTML}{17181C}   % a toolbar button of the presenter
\definecolor{guideink}{HTML}{EEF0F3}
\definecolor{guidegold}{HTML}{FFC845}   % the selected tool
\definecolor{guidepanel}{HTML}{E8EBF3}  % the slide in the gesture pictures
\definecolor{guiderule}{HTML}{C9CEDC}
\definecolor{guidemute}{HTML}{555B6B}
% guide-icons.tex -- written by make-icons.py; do not edit.
\expandafter\def\csname guideicon@open\endcsname{\draw svg {M 3 7 a 2 2 0 0 1 2 -2 h 4 l 2 2 h 8 a 2 2 0 0 1 2 2 v 8 a 2 2 0 0 1 -2 2 H 5 a 2 2 0 0 1 -2 -2 z};}
\expandafter\def\csname guideicon@reload\endcsname{\draw svg {M 19.933 13.041 a 8 8 0 1 1 -9.925 -8.788 c 3.899 -1 7.935 1.007 9.425 4.747}; \draw svg {M 20 4 v 5 h -5};}
\expandafter\def\csname guideicon@first\endcsname{\draw svg {M 11 6 l -6 6 6 6 M 19 6 l -6 6 6 6};}
\expandafter\def\csname guideicon@last\endcsname{\draw svg {M 13 6 l 6 6 -6 6 M 5 6 l 6 6 -6 6};}
\expandafter\def\csname guideicon@prev\endcsname{\draw svg {M 15 5 l -7 7 7 7};}
\expandafter\def\csname guideicon@next\endcsname{\draw svg {M 9 5 l 7 7 -7 7};}
\expandafter\def\csname guideicon@grid\endcsname{\draw[rounded corners=1pt] (3pt,5pt) rectangle (11pt,11pt); \draw[rounded corners=1pt] (13pt,5pt) rectangle (21pt,11pt); \draw[rounded corners=1pt] (3pt,13pt) rectangle (11pt,19pt); \draw[rounded corners=1pt] (13pt,13pt) rectangle (21pt,19pt);}
\expandafter\def\csname guideicon@select\endcsname{\draw svg {M 4 8 v -2 a 2 2 0 0 1 2 -2 h 2}; \draw svg {M 4 16 v 2 a 2 2 0 0 0 2 2 h 2}; \draw svg {M 16 4 h 2 a 2 2 0 0 1 2 2 v 2}; \draw svg {M 16 20 h 2 a 2 2 0 0 0 2 -2 v -2}; \draw svg {M 12 16 v -7}; \draw svg {M 9 9 h 6};}
\expandafter\def\csname guideicon@pen\endcsname{\draw svg {M 4 20 l 1.2 -4.4 L 15.6 5.2 a 2 2 0 0 1 2.8 0 l 0.4 0.4 a 2 2 0 0 1 0 2.8 L 8.4 18.8 z}; \draw svg {M 13.5 7.3 l 3.2 3.2};}
\expandafter\def\csname guideicon@highlighter\endcsname{\draw svg {M 9 14 l -4 4 v 2 h 5 l 2 -2}; \draw svg {M 9 14 l 7 -9 4 3 -7 9 z}; \draw svg {M 4 21 h 16};}
\expandafter\def\csname guideicon@laser\endcsname{\draw (12pt,12pt) circle [radius=3pt]; \draw svg {M 12 3 v 3 M 12 18 v 3 M 3 12 h 3 M 18 12 h 3};}
\expandafter\def\csname guideicon@eraser\endcsname{\draw svg {M 9 20 h 11}; \draw svg {M 4.6 14.6 l 9 -9 a 2 2 0 0 1 2.8 0 l 2 2 a 2 2 0 0 1 0 2.8 L 11 18 H 8 z};}
\expandafter\def\csname guideicon@undo\endcsname{\draw svg {M 9 6 L 4 11 l 5 5}; \draw svg {M 4 11 h 10 a 5 5 0 0 1 0 10 h -3};}
\expandafter\def\csname guideicon@clear\endcsname{\draw svg {M 4 7 h 16 M 9 7 V 4 h 6 v 3 M 6 7 l 1 13 h 10 l 1 -13};}
\expandafter\def\csname guideicon@save\endcsname{\draw svg {M 12 4 v 11 M 7 10 l 5 5 5 -5 M 5 20 h 14};}
\expandafter\def\csname guideicon@fs\endcsname{\draw svg {M 4 9 V 4 h 5 M 20 9 V 4 h -5 M 4 15 v 5 h 5 M 20 15 v 5 h -5};}
\expandafter\def\csname guideicon@hide\endcsname{\draw svg {M 6 15 l 6 -6 6 6};}
\expandafter\def\csname guideicon@hand\endcsname{\draw svg {M 7.904 17.563 a 1.2 1.2 0 0 0 2.228 0.308 l 2.09 -3.093 l 4.907 4.907 a 1.067 1.067 0 0 0 1.509 0 l 1.047 -1.047 a 1.067 1.067 0 0 0 0 -1.509 l -4.907 -4.907 l 3.113 -2.09 a 1.2 1.2 0 0 0 -0.309 -2.228 l -13.582 -3.904 l 3.904 13.563};}
\expandafter\def\csname guideicon@mouse\endcsname{\draw svg {M 6 7 a 4 4 0 0 1 4 -4 h 4 a 4 4 0 0 1 4 4 v 10 a 4 4 0 0 1 -4 4 h -4 a 4 4 0 0 1 -4 -4 l 0 -10}; \draw svg {M 12 7 l 0 4};}
\expandafter\def\csname guideicon@split\endcsname{\draw[rounded corners=1pt] (2pt,6pt) rectangle (13pt,14pt); \draw[rounded corners=1pt] (11pt,3pt) rectangle (22pt,11pt); \draw svg {M 5 19 h 6 M 8 14 v 5 M 14 11 v 4 h 5};}
\expandafter\def\csname guideicon@settings\endcsname{\draw svg {M 10.325 4.317 c 0.426 -1.756 2.924 -1.756 3.35 0 a 1.724 1.724 0 0 0 2.573 1.066 c 1.543 -0.94 3.31 0.826 2.37 2.37 a 1.724 1.724 0 0 0 1.065 2.572 c 1.756 0.426 1.756 2.924 0 3.35 a 1.724 1.724 0 0 0 -1.066 2.573 c 0.94 1.543 -0.826 3.31 -2.37 2.37 a 1.724 1.724 0 0 0 -2.572 1.065 c -0.426 1.756 -2.924 1.756 -3.35 0 a 1.724 1.724 0 0 0 -2.573 -1.066 c -1.543 0.94 -3.31 -0.826 -2.37 -2.37 a 1.724 1.724 0 0 0 -1.065 -2.572 c -1.756 -0.426 -1.756 -2.924 0 -3.35 a 1.724 1.724 0 0 0 1.066 -2.573 c -0.94 -1.543 0.826 -3.31 2.37 -2.37 c 1 0.608 2.296 0.07 2.572 -1.065}; \draw svg {M 9 12 a 3 3 0 1 0 6 0 a 3 3 0 0 0 -6 0};}
\expandafter\def\csname guideicon@screen\endcsname{\draw svg {M 3 5 a 1 1 0 0 1 1 -1 h 16 a 1 1 0 0 1 1 1 v 10 a 1 1 0 0 1 -1 1 h -16 a 1 1 0 0 1 -1 -1 v -10}; \draw svg {M 7 20 h 10}; \draw svg {M 9 16 v 4}; \draw svg {M 15 16 v 4};}
\expandafter\def\csname guideicon@info\endcsname{\draw (12pt,12pt) circle [radius=9pt]; \draw svg {M 12 8 h 0.01}; \draw svg {M 11 12 h 1 v 4 h 1};}
\expandafter\def\csname guideicon@play\endcsname{\draw svg {M 7 4 v 16 l 13 -8 z};}
\expandafter\def\csname guideicon@frames\endcsname{\draw[rounded corners=1.5pt] (3pt,7pt) rectangle (16pt,17pt); \draw svg {M 19 8.5 v 7 M 22 10 v 4};}
\expandafter\def\csname guideicon@sliders\endcsname{\draw (14pt,6pt) circle [radius=2pt]; \draw svg {M 4 6 h 8 M 16 6 h 4}; \draw (8pt,12pt) circle [radius=2pt]; \draw svg {M 4 12 h 2 M 10 12 h 10}; \draw (17pt,18pt) circle [radius=2pt]; \draw svg {M 4 18 h 11 M 19 18 h 1};}
\expandafter\def\csname guideicon@tap\endcsname{\draw svg {M 8 13 v -8.5 a 1.5 1.5 0 0 1 3 0 v 7.5}; \draw svg {M 11 11.5 v -2 a 1.5 1.5 0 0 1 3 0 v 2.5}; \draw svg {M 14 10.5 a 1.5 1.5 0 0 1 3 0 v 1.5}; \draw svg {M 17 11.5 a 1.5 1.5 0 0 1 3 0 v 4.5 a 6 6 0 0 1 -6 6 h -2 h 0.208 a 6 6 0 0 1 -5.012 -2.7 l -0.196 -0.3 c -0.312 -0.479 -1.407 -2.388 -3.286 -5.728 a 1.5 1.5 0 0 1 0.536 -2.022 a 1.867 1.867 0 0 1 2.28 0.28 l 1.47 1.47}; \draw svg {M 5 3 l -1 -1}; \draw svg {M 4 7 h -1}; \draw svg {M 14 3 l 1 -1}; \draw svg {M 15 6 h 1};}
\expandafter\def\csname guideicon@handonly\endcsname{\draw svg {M 8 13 v -8.5 a 1.5 1.5 0 0 1 3 0 v 7.5}; \draw svg {M 11 11.5 v -2 a 1.5 1.5 0 0 1 3 0 v 2.5}; \draw svg {M 14 10.5 a 1.5 1.5 0 0 1 3 0 v 1.5}; \draw svg {M 17 11.5 a 1.5 1.5 0 0 1 3 0 v 4.5 a 6 6 0 0 1 -6 6 h -2 h 0.208 a 6 6 0 0 1 -5.012 -2.7 l -0.196 -0.3 c -0.312 -0.479 -1.407 -2.388 -3.286 -5.728 a 1.5 1.5 0 0 1 0.536 -2.022 a 1.867 1.867 0 0 1 2.28 0.28 l 1.47 1.47};}

% the content area of a slide: x 0.75 to 33.11 cm, y 2.42 to 16.73 cm
\def\GX{0.75}\def\GY{2.42}\def\GW{32.36}\def\GH{14.31}

% \guideglyph[colour]{name}{x}{y}{size}{line}: the 24-unit icon with its top-left corner at (x,y),
% <size> cm wide; <line> is the line width in units of the icon (the toolbar uses 1.8). In a bdltdiagram.
\newcommand\guideglyph[6][black]{%
  \begin{scope}[shift={(#3,#4)}, scale=#5*28.4528/24, yscale=-1, draw=#1, line cap=round, line join=round,
                line width={#6*#5*28.4528/24*1pt}]
    \csname guideicon@#2\endcsname
  \end{scope}}
% \guidetile[active]{name}{x}{y}{size}: the icon on its dark tile, as a button of the toolbar
\newcommand\guidetile[5][]{%
  \def\guidetmp{#1}%
  \edef\guidefillc{\ifx\guidetmp\empty guidetile\else guidegold\fi}\edef\guideinkc{\ifx\guidetmp\empty guideink\else black\fi}%
  \fill[fill=\guidefillc, rounded corners=#5*0.25cm] (#3,#4) rectangle ++(#5,#5);
  \guideglyph[\guideinkc]{#2}{#3+#5*14/64}{#4+#5*14/64}{#5*36/64}{1.8}}
% \guidecmd: the Command key sign, drawn (Source Sans has no such character); the slides write "Ctrl/<sign>"
\newcommand\guidecmd{\tikz[baseline=-11.4pt, scale=0.6, yscale=-1, line cap=round, line join=round, line width=1.15pt]{%
  \draw svg {M 7 9 a 2 2 0 1 1 2 -2 v 10 a 2 2 0 1 1 -2 -2 h 10 a 2 2 0 1 1 -2 2 v -10 a 2 2 0 1 1 2 2 h -10};}}
% \guideicon[colour]{name}{size}: the icon as a picture of its own (for the tiles of a layout)
\newcommand\guideicon[3][white]{%
  \begin{tikzpicture}[x=1cm, y=-1cm]
    \path[use as bounding box] (0,0) rectangle (#3,#3);
    \guideglyph[#1]{#2}{0}{0}{#3}{1.5}
  \end{tikzpicture}}
% \guiderow{y}{row height}{tile size}{x}{text width}{Name}{What it does}: the words next to a tile
\newcommand\guidewords[5]{%
  \node[anchor=west, align=left, inner sep=0pt] at (#1,#2)
    {\parbox{#3cm}{\raggedright{\bdltfont{22}{100}\bfseries #4}\\[0.12cm]{\bdltfont{17}{100}\color{guidemute}#5}}};}
\newcommand\guidedot[4][]{%   [outline]{x}{y}{colour}: a colour of a tool, 1.33 cm across
  \def\guidetmp{#1}%
  \ifx\guidetmp\empty \fill[fill=#4] (#2+0.665,#3) circle [radius=0.665cm];
  \else \filldraw[fill=#4, draw=guiderule, line width=1pt] (#2+0.665,#3) circle [radius=0.665cm];\fi}
\definecolor{dotblue}{HTML}{0028A5}\definecolor{dotorange}{HTML}{FC4C02}\definecolor{dotgold}{HTML}{FFC845}
\definecolor{dotapple}{HTML}{A4D233}\definecolor{dotcyan}{HTML}{4AC9E3}\definecolor{dotberry}{HTML}{BF0D3E}
\definecolor{dotgreen}{HTML}{4E6B12}\definecolor{dotblack}{HTML}{000000}

% the rows of the tools slide: \toolrow{tile}{active}{Name}{words}{the colours of the tool}
\newcount\guiderowno
\ifguideviewer \def\ToolRH{\GH/5}\else \def\ToolRH{\GH/6}\fi
\newcommand\toolrow[5]{%
  \pgfmathsetmacro\ry{\GY+\the\guiderowno*\ToolRH+\ToolRH/2}
  \guidetile[#2]{#1}{\GX}{\ry-0.89}{1.78}
  \guidewords{\GX+2.61}{\ry-0.05}{14.4}{#3}{#4}
  \def\dx{\GX+18.06}
  \foreach \c [count=\k from 0] in {#5} {%
    \edef\cname{\c}\def\cwhite{white}%
    \ifx\cname\cwhite \guidedot[o]{\dx+\k*1.80}{\ry}{white}\else \guidedot{\dx+\k*1.80}{\ry}{dot\c}\fi}
  \ifnum\guiderowno>0 \draw[guiderule, line width=0.6pt] (\GX,\ry-\ToolRH/2) -- ++(\GW,0);\fi
  \global\advance\guiderowno by 1 }
% the items of the toolbar slide, in two columns: \baritem{icon}{second icon or empty}{Name}{words}
\newcount\guideitemno
\ifguideviewer \def\BarRH{\GH/3}\def\BarTS{1.78}\def\BarPER{3}\else \def\BarRH{\GH/4}\def\BarTS{1.78}\def\BarPER{4}\fi
\newcommand\baritem[4]{%
  \pgfmathsetmacro\cx{\GX+int(\the\guideitemno/\BarPER)*\GW/2}
  \pgfmathsetmacro\ry{\GY+mod(\the\guideitemno,\BarPER)*\BarRH+\BarRH/2}
  \guidetile{#1}{\cx}{\ry-\BarTS/2}{\BarTS}
  \def\guidesecond{#2}\ifx\guidesecond\empty\else \guidetile{#2}{\cx+\BarTS+0.25}{\ry-\BarTS/2}{\BarTS}\fi
  \guidewords{\cx+2*\BarTS+1.03}{\ry-0.03}{11.6}{#3}{#4}
  \global\advance\guideitemno by 1 }

% the slide on presenting: two blue columns, \ViewW wide with \ViewG between them, each with its
% title at the top and numbered steps below.
% \guideviewcol{column}{Title}: the block and its title; \guideviewstep{column}{row}{icon}{Name}{words}: one step
\def\ViewW{15.88}\def\ViewG{0.6}
\newcommand\guideviewcol[2]{%
  \fill[uzh@blue] (\GX+#1*\ViewW+#1*\ViewG,\GY) rectangle ++(\ViewW,\GH);
  \node[anchor=north west, text=white, inner sep=0pt] at (\GX+#1*\ViewW+#1*\ViewG+1.2,\GY+1.0) {\bdltfont{28}{100}\bfseries #2};}
\newcommand\guideviewstep[5]{%
  \guideglyph[white]{#3}{\GX+#1*\ViewW+#1*\ViewG+1.2}{\GY+3.05+#2*3.75}{2.2}{1.6}
  \node[anchor=north west, text=white, align=left, inner sep=0pt] at (\GX+#1*\ViewW+#1*\ViewG+4.1,\GY+3.0+#2*3.75)
    {\parbox{11.0cm}{\raggedright{\bdltfont{20}{100}\bfseries #4\par}\vspace{0.12cm}{\bdltfont{17}{115}#5\par}}};}

\begin{document}

% ---------------------------------------------------------------- cover
% @slide S01
\begin{frame}
  \titlepage
\end{frame}

% ---------------------------------------------------------------- three ways to use it
% @slide S02
\begin{frame}{Three ways to use it}
  \bdlticontiles[height=7.8, icon=0.6, text=22]{
    \bdlttile{\guideicon{pen}{4.4}}{With a pencil}{The pencil draws and fingers navigate. Pencil mode starts at the first touch of the pencil.}
    \bdlttile{\guideicon{mouse}{4.4}}{With a mouse}{Click to turn pages, pick a tool to draw, and press Esc to go back to Navigate.}
    \bdlttile{\guideicon{tap}{4.4}}{With a finger}{A finger works like the mouse, and two fingers always zoom.}
  }
\end{frame}

% ---------------------------------------------------------------- live: still here, alive there (the third slide: why the presenter exists, before how to use it)
% @slide S15
\begin{frame}
  \bdltfullimage{Slides that run}
    {\guideposter{media/alive/alive-poster.tex}{1147}{445}%
     \bdltwidget[id=alive, x-figure=alive, poster=blank-poster.pdf, title={A chain of blocks: a picture on the left, running on the right}]{figures/guide-figures.html}}
    {Left: what every PDF viewer shows. Right: the same figure, run by \ifguideviewer BDLT PDF Viewer\else BDLT PDF Presenter\fi.
     If pressing play changes nothing, you are reading this PDF somewhere else.}
\end{frame}

% ---------------------------------------------------------------- live media (the viewer's guide names the kinds)
\ifguideviewer
% @slide S16
\begin{frame}{Slides can be alive}
  \bdlticontiles[height=7.8, icon=0.6, text=22]{
    \bdlttile{\guideicon{play}{4.4}}{Videos and sound}{Tap the picture to play or pause, and drag the bar to move in time.}
    \bdlttile{\guideicon{frames}{4.4}}{Animations}{Let them run, or step through the frames one by one.}
    \bdlttile{\guideicon{sliders}{4.4}}{Interactive figures}{Move the sliders and drag the points with a finger or the mouse.}
  }
\end{frame}
\fi

% ---------------------------------------------------------------- the four tap zones
% @slide S06
% live: the page of a slide with its four tap zones. Each answers a tap; this is the only slide on the gestures.
\begin{frame}{The four tap zones of a slide: try them here}
  \bdltarea{0.75}{2.42}{32.36}{14.31}
    {\guideposter{media/gestures/gestures-poster.tex}{1147}{507}%
     \bdltwidget[id=gestures, x-figure=gestures, poster=blank-poster.pdf, title={This page with its four tap zones}]{figures/guide-figures.html}}
\end{frame}

% ---------------------------------------------------------------- pinch
% @slide S07
\begin{frame}{Pinch to zoom}
  \begin{bdltdiagram}
    \fill[uzh@blue] (\GX,\GY) rectangle ++(\GW,\GH);
    % two hands pointing at each other along a diagonal, arrows spreading outward (a 48-unit drawing, 10 cm)
    \begin{scope}[shift={(\GX+2.5,\GY+2.15)}, scale=10*28.4528/48, yscale=-1, draw=white, line cap=round, line join=round, line width=1.5*10*28.4528/48*1pt]
      \begin{scope}[shift={(2pt,22pt)}, rotate around={45:(12pt,12pt)}]\csname guideicon@handonly\endcsname\end{scope}
      \begin{scope}[shift={(22pt,2pt)}, rotate around={225:(12pt,12pt)}]\csname guideicon@handonly\endcsname\end{scope}
      \draw svg {M 7 19 L 19 7 M 7 19 h 4.5 M 7 19 v -4.5 M 19 7 h -4.5 M 19 7 v 4.5};
    \end{scope}
    \node[anchor=west, text=white, align=left, inner sep=0pt] at (\GX+14.5,\GY+\GH/2)
      {\parbox{17cm}{\raggedright{\bdltfont{28}{100}\bfseries Pinch with two fingers to zoom}\\[0.35cm]
        {\bdltfont{20}{130}Drag with one finger to move around while zoomed in.\\ Double-tap to fit the whole slide again.\\
         The Pencil keeps writing while you are zoomed in.}}};
  \end{bdltdiagram}
\end{frame}

% ---------------------------------------------------------------- the tools
% @slide S08
\begin{frame}{\ifguideviewer Tools: navigate, select, write, erase\else Tools: navigate, select, write, point, erase\fi}
  \begin{bdltdiagram}
    \guiderowno=0
    \toolrow{hand}{a}{Navigate}{Tap or click to turn pages (Esc returns here).}{}
    \toolrow{select}{}{Select}{Drag to copy words, or an area as a picture.}{}
    \toolrow{pen}{}{Pen}{Pressure-sensitive ink. Tap again for the colours.}{blue,orange,gold,apple,black,white}
    \toolrow{highlighter}{}{Highlighter}{Transparent marks. Tap again for the colours.}{gold,apple,cyan,orange,berry}
    \ifguideviewer\else
    \toolrow{laser}{}{Laser pointer}{Points without marking. Tap again for the colour.}{berry,green}
    \fi
    \toolrow{eraser}{}{Eraser}{Touch a stroke to remove it.}{}
  \end{bdltdiagram}
\end{frame}

% ---------------------------------------------------------------- save (after the tools: this is where the notes are kept)
% @slide S12
\begin{frame}{Save keeps your notes}
  \begin{bdltdiagram}
    \fill[uzh@blue] (\GX,\GY) rectangle ++(\GW,\GH);
    \guideglyph[white]{save}{\GX+3.05}{\GY+3.43}{7.44}{1.6}
    \node[anchor=west, text=white, align=left, inner sep=0pt] at (\GX+13.0,\GY+\GH/2)
      {\parbox{18cm}{\raggedright{\bdltfont{28}{100}\bfseries Save writes your ink into a new PDF}\\[0.35cm]
        {\bdltfont{20}{130}It works on every device: the new PDF is downloaded, or your device asks where to keep it.\\ Try it: circle something on this slide, then press Save.}}};
  \end{bdltdiagram}
\end{frame}

% ---------------------------------------------------------------- copy and paste
% @slide S17
\begin{frame}{Select: copy words or an area, and paste them anywhere}
  \bdlticontiles[height=7.8, icon=0.6, text=22]{
    \bdlttile{\guideicon{select}{4.4}}{1. Choose Select}{The button next to Navigate, or the S key. It works with the mouse; Esc goes back to Navigate.}
    \bdlttile{\guideicon{textcursor}{4.4}}{2. Copy words}{Drag over the words of a slide, then copy and paste them as text, as in any document.}
    \bdlttile{\guideicon{marquee}{4.4}}{3. Copy a picture}{Drag next to the words, or hold Alt anywhere: that area is copied as a picture, with your ink. Paste it anywhere.}
  }
\end{frame}

% ---------------------------------------------------------------- the toolbar
% @slide S09
% only the buttons that no other slide explains: the tools, Save, the views, the page turns, the
% overview and hiding the toolbar have their own slides
\begin{frame}{The other buttons of the toolbar}
  \begin{bdltdiagram}
    \guideitemno=0
    \baritem{open}{}{Open}{Ctrl/\guidecmd\,O opens a PDF, Ctrl/\guidecmd\,W closes it.}
    \ifguideviewer\else
    \baritem{reload}{}{Reload}{Reads the PDF again; the ink stays (Ctrl/\guidecmd\,R).}
    \fi
    \baritem{first}{last}{First / last slide}{Also at both ends of the slide overview.}
    \baritem{undo}{}{Undo}{Takes back the last change (Ctrl/\guidecmd\,Z).}
    \baritem{clear}{}{Clear}{Removes all ink on this slide.}
    \baritem{fs}{}{Full screen}{Hides the browser bars where allowed.}
    \baritem{settings}{}{Settings}{\ifguideviewer Colours, Pencil mode, About.\else Colours, slide strip speed and glide, About.\fi}
  \end{bdltdiagram}
\end{frame}

\ifguideviewer\else
% ---------------------------------------------------------------- presenting: one screen or two
% @slide S10
\begin{frame}{Present with the same view, or with separate views}
  \begin{bdltdiagram}
    % left: same view (one screen, mirrored); right: separate views (projector and lecturer)
    \guideviewcol{0}{Same view}
    \guideviewstep{0}{0}{mirror}{1. Mirror your screen}{Show it on the projector: everyone sees what you see.}
    \guideviewstep{0}{1}{hide}{2. Hide the toolbar}{Tap the top centre of the slide; tap there again to bring it back.}
    \guideviewcol{1}{Separate views}
    \guideviewstep{1}{0}{split}{1. Two-screens button}{Tap it in the toolbar. It says so if your device cannot open a second window.}
    \guideviewstep{1}{1}{screen}{2. Projector window}{Open it as a window of its own and move it to the projector\textquotesingle s display.}
    \guideviewstep{1}{2}{fs}{3. Fill the screen}{Tap the projector window. Your own screen then shows the clock and the slide strip.}
  \end{bdltdiagram}
\end{frame}
\fi

% @slide S13
\infoframe{Open your own PDF from the toolbar to begin}

\end{document}
